Underwater biological object detection is a practical requirement for marine resource
survey, ecological monitoring, and autonomous underwater robotics. Compared with common
in-air visual detection, underwater scenes are more challenging because of absorption,
backscatter, low contrast, color cast, and motion blur. These degradations are especially
harmful for small objects whose texture and boundary cues are already weak.

Existing studies typically improve underwater detection in two directions. The first
direction enhances detection backbones, necks, or heads to improve representation quality.
The second direction performs underwater image enhancement before detection. However, the
relationship between enhancement and detection is not always positive. Enhanced images may
recover contrast and local structure, but they can also introduce artifacts or alter the
original semantic distribution.

This paper defines underwater detection in low-visibility scenes as a cross-domain
representation problem between raw and enhanced inputs. Raw images preserve more realistic
semantic statistics, whereas enhanced images may expose details that benefit small-object
localization. Based on this observation, we design QDCR-Net, a quality-aware dual-branch
cross-residual framework that enables complementary modeling between the two domains.

The main contributions are summarized as follows:
\begin{enumerate}[leftmargin=*]
  \item We propose a dual-branch underwater detection framework that explicitly models
  raw-image semantics and enhancement-guided structure for low-visibility scenes.
  \item We design a bidirectional cross-residual interaction module to exchange stable
  semantic priors and structural detail responses across branches.
  \item We introduce a quality-aware dynamic fusion mechanism that adaptively reweights
  branch contributions according to image quality.
  \item We add a small-object-oriented neck path to preserve high-resolution detail cues
  without heavily increasing model complexity.
\end{enumerate}
